\chapter*{Introduction} \label{I.0}
\markboth{INTRODUCTION}{INTRODUCTION}

We\pageoriginale present here, in revised and expanded form, a photo-offset
reissue of the second \sisi{ALGEBRAIC GEOMETRY SEMINAR of the INSTITUT
DES HAUTES \'ETUDES SCIENTIFIQUES}{Algebraic Geometry Seminar of the
Institut des Hautes \'Etudes Scientifiques}, held in 1962 and originally
issued in mimeographed form.

The reader is referred to the Introduction to the first of these seminars
(cited below as \SGA 1) for the aims pursued by these seminars and their
relation to the \sisi{\'EL\'EMENTS DE G\'{E}OM\'ETRIE
ALG\'EBRIQUE}{\'El\'ements de G\'eom\'etrie Alg\'ebrique}.

The text of Expos\'es~\Ref{I} through~\Ref{XI} was written up as the
seminar proceeded, from my oral lectures and handwritten notes, by a group
of participants comprising I\ptbl~\sisi{GIROGUITTI}{Giorgiutti},
J\ptbl~\sisi{GIRAUD}{Giraud}, Mlle M\ptbl~\sisi{JAFFE}{Jaffe} (now
Mme M\ptbl~\sisi{HAKIM}{Hakim}), and A\ptbl~\sisi{LAUDAL}{Laudal}.
These notes were originally regarded as provisional and intended for very
limited circulation, pending their absorption into the EGA---an absorption
that has now become problematic, to say the least, as it has for the other
parts of the SGA. As the warning to the original edition said, the
\og confidential\fg character of the notes was meant to excuse certain
\og stylistic weaknesses\fg, doubtless more apparent in the present
seminar, \SGA 2, than in the others. I have tried, as far as possible, to
remedy them in the present reissue by a fairly thorough revision of the
original text. In particular, I have harmonized the systems used to number
statements in the different Expos\'es, introducing everywhere the same
decimal system, already used in most of the original Expos\'es in the
present \SGA 2 as well as in all the other parts of the SGA. This led me,
in particular, to revise completely the numbering of
the statements in Expos\'es~\Ref{III} through~\Ref{VIII}\pageoriginale
(and consequently the references to those Expos\'es)\sfootnote{It goes
without saying that every reference to \SGA 2 appearing in those parts of
the SGA published in the \sisi{SERIES IN PURE MATHEMATICS}{Series in Pure
Mathematics} refers to the present volume, not to the original edition of
\SGA 2!}. I have also tried to remove from the original text the principal
typing and syntactical errors, which were numerous and troublesome.
Moreover, Mme M\ptbl~\sisi{HAKIM}{Hakim} kindly undertook to rewrite
Expos\'e~\Ref{IV} in a less telegraphic style than the original. As in the
other SGA reissues, I have also added a number of footnotes, either to give
additional references or to indicate the status of questions on which
progress has been made since the original text was written. Finally, the
seminar has been enlarged by a new Expos\'e, \sisi{namely}{namely}
Expos\'e~\Ref{XIV}, written by Mme\sisi{.}{}
\sisi{MICH\`ELE RAYNAUD}{Mich\`ele Raynaud} in 1967, which takes up and
completes suggestions contained in the \og Comments on
Expos\'e~\Ref{XIII}\fg (\Ref{XIII}~\Ref{XIII.6}), written in March 1963.
This Expos\'e treats Lefschetz-type theorems from the viewpoint of \'etale
cohomology, using the results on \'etale cohomology presented in \SGA 4
and \SGA 5 (to appear in the same
\sisi{SERIES IN PURE MATHEMATICS}{Series in Pure Mathematics}); it is therefore less
\og elementary\fg in nature than the other Expos\'es of the present
volume, which use little more than the substance of Chapters~I through~III
of the EGA. Here is an outline of the contents of the present volume.
Expos\'e~\Ref{I} contains the sorites of \og cohomology with support in
$Y$\fg, $H_Y^*(X,F)$, where $Y$ is a closed subset of a space $X$,
cohomology that may be interpreted as the cohomology of $X$ \emph{modulo}
the open subset $X-Y$, and that is the abutment of a very useful
\og \emph{local-to-global spectral sequence}\fg
\Ref{I}~\Ref{I.2.6}, involving the \emph{sheaves} of cohomology
\og with support in $Y$\fg, $\sheaf{H}_Y^*(F)$. In many questions this
formalism can\pageoriginale play a \og localization\fg role analogous to
that played by tubular neighborhoods of $Y$ in differential geometry.
Expos\'e~\Ref{II} studies the preceding notions for quasi-coherent sheaves
on preschemes, and Expos\'e~\Ref{III} relates them to the classical notion
of \emph{depth} (\Ref{III}~\Ref{III.3.3}).

Expos\'es~\Ref{IV} and~\Ref{V} give forms of \emph{local duality}, which
may be compared with Serre's projective duality theorem
(\Ref{XII}~\Ref{XII.1.1}). We note that both kinds of duality theorem are
substantially generalized in \sisi{HARTSHORNE}{Hartshorne}'s seminar
(cited in the footnote at the end of \Exp \Ref{IV})\refstepcounter{toto}

Expos\'es~\Ref{VI} and~\Ref{VII} give elementary technical notions used in
Expos\'e~\Ref{VIII} to prove the \emph{finiteness theorem}
(\Ref{VIII}~\Ref{VIII.2.3}). This theorem gives necessary and sufficient
conditions, for a coherent sheaf $F$ on a noetherian scheme $X$, for the
local cohomology sheaves $\sheaf{H}_Y^i(F)$ to be coherent for $i\leq n$
(or, equivalently, for the sheaves $\R^if_*(F|X-Y)$ to be coherent for
$i\leq n-1$, where $f\colon X-Y\to X$ is the inclusion). It is one of the
central technical results of the seminar, and in Expos\'e~\Ref{IX} we show
how a theorem of this kind can be used to establish a \og comparison
theorem\fg and an \og existence theorem\fg in formal geometry, following
and generalizing the use made in (\EGA III~\oldS\S 4 and~5) of the
finiteness theorem for a proper morphism.

These last results are applied in~\Ref{X} and~\Ref{XI}, devoted
respectively to Lefschetz-type theorems for the fundamental group and
for\pageoriginale the Picard group. These theorems compare, under certain
conditions, the invariants ($\pi_1$ or $\Pic$) associated respectively
with a scheme $X$ and a subscheme $Y$, which plays the role of a
hyperplane section, and in particular give conditions under which they
are isomorphic. Roughly speaking, the hypotheses allow one to pass from
$Y$ to the formal completion of $X$ along $Y$, and then to apply the
results of~\Ref{IX} in order to pass from there to an open
\emph{neighborhood} $U$ of $Y$ in $X$. To pass from $U$ to $X$, one also
needs information of a \og purity\fg or \og parafactoriality\fg type for
the local rings of $X$ at the points of $Z=X-U$, which is a finite
discrete set in the cases under consideration. This explains the
interaction, in the proofs of Expos\'es~\Ref{X}, \Ref{XI}, and \Ref{XII},
between local and global results, particularly in certain inductions. The
main results obtained in~\Ref{X} and~\Ref{XI} are the theorems of
\emph{local nature} \Ref{X}~\Ref{X.3.4} (the \emph{purity theorem}) and
\Ref{XI}~\Ref{XI.3.14} (the \emph{parafactoriality theorem}). Notice that
these theorems are proved by cohomological techniques of an essentially
global nature. In~\Ref{XII}, using the preceding local results, one
obtains global variants for projective schemes over a field, or more
generally over a more or less arbitrary base scheme; typical statements
include \Ref{XII}~\Ref{XII.3.5} and \Ref{XII}~\Ref{XII.3.7}.

In~\Ref{XIII}, we review some of the many problems and conjectures
suggested by the results and methods of the seminar. Perhaps the most
interesting concern cohomological and homotopical Lefschetz-type theorems
for complex analytic spaces; \cf \Ref{XIII}, pages 26 and following. In the context of the \'etale cohomology of
schemes\pageoriginale, the corresponding conjectures are proved
in~\Ref{XIV} by a duality technique that should apply equally in the
complex analytic case (\cf the comments in \Ref{XIII}, p\ptbl~25, and
\Ref{XIV}~\Ref{XIV.6.4}). But the corresponding homotopical statements
for analytic spaces, and more particularly those involving the
fundamental group, seem to require entirely new techniques
(\cf \Ref{XIV}~\Ref{XIV.6.4}).

I am happy to thank all those who, in various capacities, helped bring the
present volume to publication, including the collaborators already named
in this Introduction. In particular, I wish to thank Mlle
\sisi{CHARDON}{Chardon} for the good grace with which she carried out the
thankless task of physically preparing the final manuscript for
photo-offset reproduction.

\bigskip\hfill Bures-sur-Yvette, April~1968
\par\smallskip \hfill A\ptbl~Grothendieck.
